% S4: ban TIENG VIET cua paper2/blocks/08_subsidence.tex.
% Chi duoc \input boi paper2_hujos/supplementary_vi.tex.
% NGUYEN TAC: giu nguyen moi con so va moi hedge; thap phan kieu Viet ({,});
% tham chieu cheo toi bai chinh dung \eqref/\ref (so lay tu main_hujos_vi.aux
% boi experiments/make_supp_stubs_vi.py) thay vi go tay "Eq.~(16)" — ban EN go tay
% nen da lech so voi ban VI (eq:umax EN=16, VI=15).
\section*{S4. Tính bất biến đối với sụt lún}

Tuân thủ AWD thường được phát biểu bằng một mực nước ngầm hoặc độ sâu ngập tuyệt
đối, ví dụ $15$\,cm dưới mặt ruộng, tức ngưỡng tưới lại của nghiệm thức Safe-AWD
trong~\cite{balaine2019}. Ở một đồng bằng đang sụt lún, cao độ mặt ruộng $z_s(t)$ hạ
với tốc độ $\dot z$, nên một ngưỡng tuyệt đối đo theo một mốc cố định sẽ trôi.

\begin{proposition}[Độ trôi của một ngưỡng tuyệt đối]
\label{prop:subsid}
Trong một chu kỳ cấp tín chỉ dài $T_c$ năm, một ngưỡng tuyệt đối biểu diễn theo một
cao độ gốc cố định sẽ dịch đi $\dot z\,T_c$. Với $\dot z\in[1{,}1;2{,}5]$~cm/năm và
$T_c=5$~năm, độ dịch là $5{,}5$--$12{,}5$~cm, tức $37$--$83\%$ của một ngưỡng
$15$\,cm. Một tiêu chí viết theo tỷ lệ cạn kiệt \eqref{eq:depl} có độ trôi bằng
không, vì $FC$ và $WP$ là tính chất của đất quy về tầng rễ chứ không quy về một cao
độ tuyệt đối.
\end{proposition}

\begin{proof}
Tiêu chí tuyệt đối so một độ sâu đo được với một mốc không đổi; nếu mốc và mặt ruộng
rời nhau với tốc độ $\dot z$ thì ngưỡng hữu hiệu thay đổi $\dot z T_c$ sau $T_c$ năm.
Tiêu chí cạn kiệt so $FC-w_t$ với $(FC-WP)$, cả hai đều được định nghĩa bên trong tầng
rễ và tịnh tiến cứng cùng mặt ruộng, nên tỷ số của chúng bất biến đối với $z_s$.
\end{proof}

Các tốc độ sụt lún là số đo cho đồng bằng này trong~\cite{minderhoud2017}. Hệ quả thực
tiễn là một phương án neo theo tỷ lệ cạn kiệt không cần đo đạc lại các mốc ruộng trong
suốt chu kỳ cấp tín chỉ năm năm, còn phương án dùng độ sâu tuyệt đối thì cần.
